\section{KẾT LUẬN VÀ HƯỚNG PHÁT TRIỂN}

\subsection{Kết luận}

Đồ án này đã trình bày một quy trình nghiên cứu toàn diện về việc xây dựng hệ thống phát hiện mã độc phân tán, kết hợp giữa sức mạnh của Federated Learning và tiềm năng của Quantum Machine Learning. Dựa trên các kết quả thực nghiệm trên bộ dữ liệu \textit{Obfuscated-MalMem2022}, chúng tôi rút ra các kết luận chính sau:

\begin{enumerate}
    \item \textbf{Hiệu quả của Federated Learning:} Hệ thống được xây dựng trên nền tảng Flower Framework \cite{beutel2022flowerfriendlyfederatedlearning} đã chứng minh khả năng hoạt động ổn định và đạt độ chính xác rất cao ($> 99.7\%$) trên cả mô hình \textbf{Logistic Regression} và \textbf{Multi-layer Perceptron (MLP)}, khẳng định tính khả thi của việc phát hiện mã độc phân tán mà không cần thu thập dữ liệu thô về máy chủ trung tâm.
    
    \item \textbf{Tiềm năng của Quantum Machine Learning:} Mặc dù bị giới hạn về tài nguyên mô phỏng, mô hình \textbf{Hybrid QNN} với lớp nén 4-qubit đạt độ chính xác kiểm thử lên đến \textbf{99.92\%} sau 5 vòng huấn luyện liên kết trong thử nghiệm sơ bộ. Kết quả cho thấy tiềm năng của kiến trúc lai ghép trong việc biểu diễn thông tin cô đọng thông qua không gian Hilbert (giảm chiều dữ liệu từ 55 đặc trưng đầu vào xuống còn 4 chiều). Tuy nhiên, cần lưu ý rằng phần lớn năng lực biểu diễn đến từ encoder cổ điển (mạng nơ-ron nén 55→4), và cần thêm các thí nghiệm có đối chứng để xác nhận chắc chắn đóng góp của lớp lượng tử.
    
    \item \textbf{Tính minh bạch và An toàn:} Việc tích hợp cơ chế bảo mật \textbf{SSL/TLS phía máy chủ (Server-side TLS)} đảm bảo kênh truyền thông tin cậy và an toàn trước các nguy cơ nghe lén. Đồng thời, module \textbf{XAI} tích hợp trên Dashboard giúp giải mã quyết định của mô hình, chỉ ra mức độ quan trọng của các đặc trưng như luồng tiến trình (\texttt{pslist.avg\_threads}), thư viện động (\texttt{modules.nmodules}) hay dịch vụ (\texttt{svcscan.process\_services}), hỗ trợ đắc lực cho các chuyên gia phân tích.
    
    \item \textbf{Lựa chọn mô hình tối ưu:} Thông qua quá trình benchmark trên tập trung, nghiên cứu khẳng định sự vượt trội của các thuật toán Gradient Boosting (CatBoost, XGBoost), nhưng cũng biện luận thuyết phục cho việc lựa chọn các lớp PyTorch (Logistic Regression, MLP, Hybrid QNN) làm nền tảng cho FL nhờ tính tương thích đạo hàm và khả năng cập nhật trọng số liên kết.
\end{enumerate}

\subsection{Hạn chế và Hướng phát triển}

Bên cạnh những kết quả đạt được, đề tài vẫn còn một số hạn chế về mặt phần cứng và quy mô thử nghiệm. Trong tương lai, nhóm nghiên cứu đề xuất các hướng mở rộng sau:

\subsubsection{Triển khai trên Phần cứng Lượng tử Thực (Real Quantum Hardware)}
Hiện tại, các thực nghiệm QML đang được chạy trên máy giả lập (Simulator) sử dụng CPU/GPU, dẫn đến thời gian huấn luyện kéo dài. Hướng phát triển tiếp theo là kết nối hệ thống với các bộ xử lý lượng tử (QPU) thực tế thông qua các dịch vụ đám mây như \textit{IBM Quantum Experience} hoặc \textit{Amazon Braket} để đánh giá tốc độ xử lý (Quantum Speedup) trong điều kiện thực tế.


\subsubsection{Mở rộng Quy mô và Chống chịu Tấn công}
Thử nghiệm hiện tại giới hạn ở số lượng client nhỏ ($K=2$). Chúng tôi dự định mở rộng hệ thống lên hàng chục hoặc hàng trăm client giả lập (sử dụng Docker/Kubernetes) để kiểm tra khả năng chịu tải. Đồng thời, nghiên cứu sâu hơn về các chiến lược phòng thủ trước các cuộc tấn công phức tạp như \textit{Backdoor Attack} hay \textit{Sybil Attack} trong môi trường FL quy mô lớn.

\subsubsection{Triển khai trên Thiết bị Biên Vật lý}
Thay vì chỉ chạy giả lập trên máy tính, hệ thống sẽ được đóng gói và triển khai thử nghiệm trên các thiết bị biên thực tế có tài nguyên hạn chế như \textit{Raspberry Pi} hoặc \textit{NVIDIA Jetson Nano}. Điều này sẽ giúp đánh giá chính xác hơn về mức tiêu thụ năng lượng và độ trễ thực tế của mô hình QNN rút gọn.